\documentclass[12pt, a4paper]{article}

% ── Shared preamble (src/lecturenotes.sty) ────────────────────────────────────
%   Packages, colors, boxes, \sol/\soleq and math macros all live in
%   src/*.sty. \input@path points LaTeX two folders up to find them.
\makeatletter
\def\input@path{{../../}}
\makeatother
\usepackage{lecturenotes}

% ── Local TikZ styles ─────────────────────────────────────────────────────────
\tikzset{
  neuron/.style={circle, draw=mainblue, fill=boxbg,
      minimum size=1.1cm, thick, font=\small},
  preact/.style={circle, draw=mainblue, fill=white,
      minimum size=1.1cm, thick, font=\small},
  input/.style={circle, draw=gray!60, fill=gray!15,
      minimum size=1.1cm, thick, font=\small},
  lossnode/.style={circle, draw=red!70!black, fill=warnbg,
      minimum size=1.1cm, thick, font=\small},
  unitbox/.style={draw=gray!60, dashed, rounded corners=6pt},
  back/.style={->, thick, red!70!black, dashed},
}

% ── Document ──────────────────────────────────────────────────────────────────
\begin{document}

\begin{center}
  {\LARGE \textbf{Section 3.3: Forward Pass and Backpropagation}} \\[0.3em]
  {\large\itshape From Prediction to Learning} \\[0.5em]
  {\large AI-IM Class --- Lecture Notes}%
  \ifstudentmode
    \quad {\normalsize\color{gray}[Student Copy]}%
  \fi
\end{center}

\vspace{1em}
\hrule
\vspace{1.5em}

In Section 3.2 we learned that gradient descent improves a model by
repeating two things: compute the loss, then nudge every parameter against
its gradient. For a neural network this cycle has two halves:
\begin{itemize}[nosep]
  \item the \textbf{forward pass}, which turns an input $\bx$ into a
        prediction $\yhat$ and a loss $J$, and
  \item \textbf{backpropagation}, which works out how much each weight
        contributed to that loss, giving the gradient $\partial J/\partial w$.
\end{itemize}
We will compute the forward pass by hand, see the \emph{idea} behind
backpropagation, and then check that one update step really lowers the loss.

% ══════════════════════════════════════════════════════════════════════════════
\section{The Network}

We work with a \textbf{two-layer feedforward network} (one hidden layer)
for the XOR problem:

\begin{figure}[H]
  \centering
  \begin{tikzpicture}[>=Stealth, node distance=0pt]
    % ── Nodes ──────────────────────────────────────────────────────────────────
    \node[input]   (x1) at (0,  1.4) {$x_1$};
    \node[input]   (x2) at (0, -1.4) {$x_2$};
    \node[neuron]  (h1) at (4,  1.4) {$h_1$};
    \node[neuron]  (h2) at (4, -1.4) {$h_2$};
    \node[neuron]  (yh) at (8,  0)   {$\yhat$};

    % ── Connections: input → hidden ────────────────────────────────────────────
    \draw[->, thick] (x1) -- node[above, font=\scriptsize, pos=0.5]
    {$w^{(1)}_{11}=\phantom{-}0.5$} (h1);
    \draw[->, thick] (x1) -- node[above left, font=\scriptsize, pos=0.35]
    {$w^{(1)}_{21}=\phantom{-}0.5$} (h2);
    \draw[->, thick] (x2) -- node[below left, font=\scriptsize, pos=0.35]
    {$w^{(1)}_{12}=-0.5$} (h1);
    \draw[->, thick] (x2) -- node[below, font=\scriptsize, pos=0.5]
    {$w^{(1)}_{22}=\phantom{-}0.5$} (h2);

    % ── Connections: hidden → output ───────────────────────────────────────────
    \draw[->, thick] (h1) -- node[above, font=\scriptsize, pos=0.5]
    {$w^{(2)}_1=0.5$} (yh);
    \draw[->, thick] (h2) -- node[below, font=\scriptsize, pos=0.5]
    {$w^{(2)}_2=0.5$} (yh);

    % ── Bias annotations ───────────────────────────────────────────────────────
    \node[font=\scriptsize, gray!80!black, above=0.25cm of h1] {$b^{(1)}_1 = 0.5$};
    \node[font=\scriptsize, gray!80!black, below=0.25cm of h2] {$b^{(1)}_2 = 0.0$};
    \node[font=\scriptsize, gray!80!black, above=0.25cm of yh] {$b^{(2)} = 0.0$};

    % ── Layer labels ───────────────────────────────────────────────────────────
    \node[font=\small\itshape, gray] at (0, 3.1) {Input layer};
    \node[font=\small\itshape, gray] at (4, 3.1) {Hidden layer ($\sigmoid$)};
    \node[font=\small\itshape, gray] at (8, 3.1) {Output ($\sigmoid$)};
  \end{tikzpicture}
  \caption{A two-layer feedforward network with one hidden layer.}
  \label{fig:architecture}
\end{figure}

In matrix form, the (not yet trained) parameters are
\[
  \bW^{(1)} = \begin{bmatrix} 0.5 & -0.5 \\ 0.5 & \phantom{-}0.5 \end{bmatrix},
  \quad
  \bb^{(1)} = \begin{bmatrix} 0.5 \\ 0.0 \end{bmatrix},
  \quad
  \bW^{(2)} = \begin{bmatrix} 0.5 & 0.5 \end{bmatrix},
  \quad
  b^{(2)} = 0.0,
\]
where row $i$ of $\bW^{(1)}$ holds the weights \emph{feeding into} hidden
unit $i$.

% ──────────────────────────────────────────────────────────────────────────────
\subsection{The Sigmoid Activation}

In Section 3.1 we used the \textbf{step function}. Its derivative is $0$
almost everywhere, so gradient descent would get no signal about which way
to move the weights. We therefore use its smooth ``S-shaped'' version, the
\textbf{sigmoid}:
\[
  \sigmoid(z) = \frac{1}{1 + e^{-z}}, \qquad \sigmoid(z) \in (0, 1).
\]

Every hidden and output neuron does two sub-steps:
\begin{enumerate}[nosep]
  \item \textbf{Pre-activation} $z$: a weighted sum of inputs plus bias.
  \item \textbf{Activation} $h = \sigmoid(z)$: squeeze $z$ into $(0,1)$.
\end{enumerate}

\begin{figure}[H]
  \centering
  \begin{tikzpicture}[>=Stealth]
    % ── Nodes ──────────────────────────────────────────────────────────────────
    \node[input]    (x1) at (0,    1.4) {$x_1$};
    \node[input]    (x2) at (0,   -1.4) {$x_2$};
    \node[preact]   (z1) at (3.6,  1.4) {$z^{(1)}_1$};
    \node[neuron]   (h1) at (5.6,  1.4) {$h_1$};
    \node[preact]   (z2) at (3.6, -1.4) {$z^{(1)}_2$};
    \node[neuron]   (h2) at (5.6, -1.4) {$h_2$};
    \node[preact]   (zo) at (9.2,  0)   {$z^{(2)}$};
    \node[neuron]   (yh) at (11.2, 0)   {$\yhat$};
    \node[lossnode] (J)  at (13.6, 0)   {$J$};
    \node[input]    (y)  at (13.6, -2.0) {$y$};

    % ── Dashed boxes: one box = one neuron ─────────────────────────────────────
    \draw[unitbox] (2.85,  0.65) rectangle (6.35,  2.15);
    \draw[unitbox] (2.85, -2.15) rectangle (6.35, -0.65);
    \draw[unitbox] (8.45, -0.75) rectangle (11.95, 0.75);

    % ── Connections ────────────────────────────────────────────────────────────
    \draw[->, thick] (x1) -- node[above, font=\scriptsize] {$w^{(1)}_{11}$} (z1);
    \draw[->, thick] (x1) -- node[sloped, above, font=\scriptsize, pos=0.28] {$w^{(1)}_{21}$} (z2);
    \draw[->, thick] (x2) -- node[sloped, below, font=\scriptsize, pos=0.28] {$w^{(1)}_{12}$} (z1);
    \draw[->, thick] (x2) -- node[below, font=\scriptsize] {$w^{(1)}_{22}$} (z2);
    \draw[->, thick, mainblue] (z1) -- node[above, font=\small] {$\sigmoid$} (h1);
    \draw[->, thick, mainblue] (z2) -- node[above, font=\small] {$\sigmoid$} (h2);
    \draw[->, thick, mainblue] (zo) -- node[above, font=\small] {$\sigmoid$} (yh);
    \draw[->, thick] (h1) -- node[above right, font=\scriptsize, pos=0.45] {$w^{(2)}_1$} (zo);
    \draw[->, thick] (h2) -- node[below right, font=\scriptsize, pos=0.45] {$w^{(2)}_2$} (zo);
    \draw[->, thick, red!70!black] (yh) -- (J);
    \draw[->, thick, red!70!black] (y)  -- (J);

    % ── Annotations ────────────────────────────────────────────────────────────
    \node[font=\scriptsize, gray!80!black] at (3.6,  2.45) {$+\,b^{(1)}_1$};
    \node[font=\scriptsize, gray!80!black] at (3.6, -2.45) {$+\,b^{(1)}_2$};
    \node[font=\scriptsize, gray!80!black] at (9.2,  1.05) {$+\,b^{(2)}$};
    \node[font=\scriptsize, red!70!black] at (13.6, 0.95) {$J = (y - \yhat)^2$};
    \node[font=\scriptsize, gray!80!black, left=0.1cm of y] {target};

    % ── Layer labels ───────────────────────────────────────────────────────────
    \node[font=\small\itshape, gray] at (0,    3.1) {Input layer};
    \node[font=\small\itshape, gray] at (4.6,  3.1) {Hidden layer};
    \node[font=\small\itshape, gray] at (10.2, 3.1) {Output layer};
    \node[font=\small\itshape, gray] at (13.6, 3.1) {Loss};
  \end{tikzpicture}
  \caption{The same network with each neuron (dashed box) split into its
    pre-activation $z$ (white) and activation (blue). The prediction $\yhat$
    is compared with the target $y$ to give the loss $J$.}
  \label{fig:forward-graph}
\end{figure}

\begin{codebox}
import numpy as np

def sigmoid(x):
    return 1 / (1 + np.exp(-x))
\end{codebox}

% ══════════════════════════════════════════════════════════════════════════════
\section{The Forward Pass}

We feed the XOR input $\bx = (x_1, x_2) = (0, 1)$, whose target is $y = 1$.

% ──────────────────────────────────────────────────────────────────────────────
\begin{stepbox}{Step 1 \quad Hidden layer}
  Weighted sums entering each hidden neuron:
  \begin{align*}
    z^{(1)}_1 & = w^{(1)}_{11}x_1 + w^{(1)}_{12}x_2 + b^{(1)}_1
    = (0.5)(0) + (-0.5)(1) + 0.5 = \sol[2.5cm]{0.0}             \\
    z^{(1)}_2 & = w^{(1)}_{21}x_1 + w^{(1)}_{22}x_2 + b^{(1)}_2
    = (0.5)(0) + (0.5)(1) + 0.0 = \sol[2.5cm]{0.5}
  \end{align*}
  Apply the sigmoid:
  \[
    h_1 = \sigmoid(z^{(1)}_1) = \sol[2.5cm]{0.500},
    \qquad
    h_2 = \sigmoid(z^{(1)}_2) = \sol[2.5cm]{0.622}
  \]
\end{stepbox}

\vspace{0.8em}

% ──────────────────────────────────────────────────────────────────────────────
\begin{stepbox}{Step 2 \quad Output layer}
  \[
    z^{(2)} = w^{(2)}_1 h_1 + w^{(2)}_2 h_2 + b^{(2)}
    = (0.5)(0.500) + (0.5)(0.622) + 0.0 = \sol[2.5cm]{0.561}
  \]
  \[
    \yhat = \sigmoid(z^{(2)}) = \sol[2.5cm]{0.637}
  \]
\end{stepbox}

\vspace{0.8em}

% ──────────────────────────────────────────────────────────────────────────────
\begin{stepbox}{Step 3 \quad Loss}
  Using the squared error:
  \[
    J = (y - \yhat)^2 = (1 - 0.637)^2 \approx \sol[2.5cm]{0.132}
  \]
\end{stepbox}

The network outputs $0.637$ when it should output $1$. A perfect prediction
would give $J = 0$. Our goal now is to change the weights so that $J$ goes
down.

% ══════════════════════════════════════════════════════════════════════════════
\section{Backpropagation: The Idea}

Gradient descent (Section 3.2) updates every parameter with
\[
  w \;\longleftarrow\; w - \eta \,\pd{J}{w}.
\]
So we need $\partial J / \partial w$ for \emph{every} weight and bias:
nine numbers for our small network, and millions for a real one.
The difficulty is that the loss depends on a weight only \emph{indirectly},
through a chain of intermediate values:
\[
  w^{(1)}_{12} \;\to\; z^{(1)}_1 \;\to\; h_1 \;\to\; z^{(2)}
  \;\to\; \yhat \;\to\; J.
\]

% ──────────────────────────────────────────────────────────────────────────────
\subsection{The Chain Rule}

If $y$ depends on $u$, and $u$ depends on $x$, then
\[
  \frac{dy}{dx} = \frac{dy}{du} \cdot \frac{du}{dx}.
\]
Longer chains simply multiply more factors.

\begin{example}
  Let $y = (2x+1)^2$. Write $u = 2x+1$ so that $y = u^2$:
  \[
    \frac{dy}{dx}
    = \underbrace{\frac{dy}{du}}_{2u}
    \cdot
    \underbrace{\frac{du}{dx}}_{2}
    = \sol[4cm]{4(2x+1)}.
  \]
\end{example}

% ──────────────────────────────────────────────────────────────────────────────
\subsection{Walking Backwards Through the Network}

Each arrow in Figure~\ref{fig:forward-graph} is a simple function whose
derivative is easy (a weighted sum or a sigmoid). The chain rule says:
\begin{center}
  \emph{the gradient of $J$ with respect to a weight is the product of the
    local derivatives along the path from $J$ back to that weight.}
\end{center}

\begin{figure}[H]
  \centering
  \begin{tikzpicture}[>=Stealth]
    % ── Nodes ──────────────────────────────────────────────────────────────────
    \node[input]    (x1) at (0,    1.4) {$x_1$};
    \node[input]    (x2) at (0,   -1.4) {$x_2$};
    \node[preact]   (z1) at (3.6,  1.4) {$z^{(1)}_1$};
    \node[neuron]   (h1) at (5.6,  1.4) {$h_1$};
    \node[preact]   (z2) at (3.6, -1.4) {$z^{(1)}_2$};
    \node[neuron]   (h2) at (5.6, -1.4) {$h_2$};
    \node[preact]   (zo) at (9.2,  0)   {$z^{(2)}$};
    \node[neuron]   (yh) at (11.2, 0)   {$\yhat$};
    \node[lossnode] (J)  at (13.6, 0)   {$J$};

    % ── Forward connections (faded) ────────────────────────────────────────────
    \foreach \a/\b in {x1/z1, x1/z2, x2/z1, x2/z2, z1/h1, z2/h2,
        h1/zo, h2/zo, zo/yh, yh/J}
    \draw[->, gray!50] (\a) -- (\b);

    % ── Backward (chain-rule) arrows ───────────────────────────────────────────
    \draw[back] (J)  to[bend right=40] node[above, font=\scriptsize] {$\pd{J}{\yhat}$} (yh);
    \draw[back] (yh) to[bend right=40] node[above, font=\scriptsize] {$\pd{\yhat}{z^{(2)}}$} (zo);
    \draw[back] (zo) to[bend right=25] node[above right, font=\scriptsize] {$\pd{z^{(2)}}{h_1}$} (h1);
    \draw[back] (zo) to[bend left=25]  node[below right, font=\scriptsize] {$\pd{z^{(2)}}{h_2}$} (h2);
    \draw[back] (h1) to[bend right=45] node[above, font=\scriptsize] {$\pd{h_1}{z^{(1)}_1}$} (z1);
    \draw[back] (h2) to[bend left=45]  node[below, font=\scriptsize] {$\pd{h_2}{z^{(1)}_2}$} (z2);
  \end{tikzpicture}
  \caption{Backpropagation runs the forward graph in reverse. Starting from
    $J$, it multiplies the local derivative on each dashed arrow. Paths share
    their first links ($J \to \yhat \to z^{(2)}$), so those products are
    computed once and reused for every weight further back.}
  \label{fig:backward-graph}
\end{figure}

For example, the gradient for $w^{(1)}_{12}$ is the product along its path:
\[
  \pd{J}{w^{(1)}_{12}}
  =
  \underbrace{\pd{J}{\yhat}}_{2(\yhat - y)}
  \cdot
  \underbrace{\pd{\yhat}{z^{(2)}}}_{\sigmoid'(z^{(2)})}
  \cdot
  \underbrace{\pd{z^{(2)}}{h_1}}_{w^{(2)}_1}
  \cdot
  \underbrace{\pd{h_1}{z^{(1)}_1}}_{\sigmoid'(z^{(1)}_1)}
  \cdot
  \underbrace{\pd{z^{(1)}_1}{w^{(1)}_{12}}}_{x_2}.
\]

\begin{patternbox}
  \textbf{What to remember about backpropagation.}
  \begin{enumerate}[nosep]
    \item It starts at the loss with the \textbf{error}
          $\pd{J}{\yhat} = 2(\yhat - y)$, i.e.\ ``prediction minus target''.
    \item Moving back one layer multiplies in that layer's weight and one
          activation slope $\sigmoid'(z)$. For the sigmoid,
          $\sigmoid'(z) = \sigmoid(z)\bigl(1-\sigmoid(z)\bigr) \le 0.25$.
    \item At the end, each weight's gradient is
          \[
            \pd{J}{w} = (\text{error signal at its neuron})
            \times (\text{input to that weight}).
          \]
          A weight whose input is $0$ gets zero gradient: it had no effect on
          this prediction.
  \end{enumerate}
\end{patternbox}

\begin{notebox}
  \textbf{Nobody does this by hand.} Deep-learning libraries record the
  forward pass and apply the chain rule automatically
  (\emph{automatic differentiation}). In PyTorch, for example, a single
  call to \texttt{loss.backward()} fills in $\partial J/\partial w$ for
  every parameter.
\end{notebox}

% ══════════════════════════════════════════════════════════════════════════════
\section{One Update Step}

Running backpropagation on our input $(0, 1)$ and applying
$w \leftarrow w - \eta\,\partial J/\partial w$ with learning rate
$\eta = 0.1$ gives the following new parameters:

\begin{center}
  \renewcommand{\arraystretch}{1.4}
  \begin{tabular}{lrrr}
    \toprule
    Parameter      & Old              & Gradient $\partial J/\partial(\cdot)$ & New                           \\
    \midrule
    $w^{(2)}_1$    & $\phantom{-}0.5$ & $-0.084$                              & $\phantom{-}0.508$            \\
    $w^{(2)}_2$    & $\phantom{-}0.5$ & $-0.105$                              & $\phantom{-}0.510$            \\
    $b^{(2)}$      & $\phantom{-}0.0$ & $-0.168$                              & $\phantom{-}0.017$            \\
    \midrule
    $w^{(1)}_{11}$ & $\phantom{-}0.5$ & $\phantom{-}0.000$                    & $\phantom{-}0.500$            \\
    $w^{(1)}_{12}$ & $-0.5$           & $-0.021$                              & $-0.498$                      \\
    $b^{(1)}_1$    & $\phantom{-}0.5$ & $-0.021$                              & $\phantom{-}0.502$            \\
    \midrule
    $w^{(1)}_{21}$ & $\phantom{-}0.5$ & $\phantom{-}0.000$                    & \sol[2cm]{$\phantom{-}0.500$} \\
    $w^{(1)}_{22}$ & $\phantom{-}0.5$ & $-0.020$                              & \sol[2cm]{$\phantom{-}0.502$} \\
    $b^{(1)}_2$    & $\phantom{-}0.0$ & $-0.020$                              & \sol[2cm]{$\phantom{-}0.002$} \\
    \bottomrule
  \end{tabular}
\end{center}

Notice that:
\begin{itemize}[nosep]
  \item All non-zero gradients are \textbf{negative}, so every parameter that
        moves goes \emph{up}, which pushes $\yhat$ up towards $y = 1$.
  \item $w^{(1)}_{11}$ and $w^{(1)}_{21}$ have zero gradient because their
        input is $x_1 = 0$, so they stay put.
  \item The output-layer gradients are larger than the hidden-layer ones:
        each step back multiplies in another factor $\sigmoid'(z) \le 0.25$.
\end{itemize}

% ──────────────────────────────────────────────────────────────────────────────
\subsection{Did the Loss Go Down?}

Run the forward pass again on the same input $(x_1, x_2) = (0, 1)$, this
time with the \textbf{new} parameters.

\begin{stepbox}{Step 1 \quad Hidden layer}
  \begin{align*}
    z^{(1)}_1 & = (0.500)(0) + (-0.498)(1) + 0.502 = \sol[2.5cm]{0.004} \\
    z^{(1)}_2 & = (0.500)(0) + (0.502)(1) + 0.002 = \sol[2.5cm]{0.504}
  \end{align*}
  \[
    h_1 = \sigmoid(z^{(1)}_1) = \sol[2.5cm]{0.501},
    \qquad
    h_2 = \sigmoid(z^{(1)}_2) = \sol[2.5cm]{0.623}
  \]
\end{stepbox}

\vspace{0.8em}

\begin{stepbox}{Step 2 \quad Output layer}
  \[
    z^{(2)} = (0.508)(0.501) + (0.510)(0.623) + 0.017 = \sol[2.5cm]{0.589}
  \]
  \[
    \yhat = \sigmoid(z^{(2)}) = \sol[2.5cm]{0.643}
  \]
\end{stepbox}

\vspace{0.8em}

\begin{stepbox}{Step 3 \quad Loss}
  \[
    J = (1 - 0.643)^2 \approx \sol[2.5cm]{0.127}
  \]
\end{stepbox}

\begin{center}
  \renewcommand{\arraystretch}{1.4}
  \begin{tabular}{lcc}
    \toprule
    Quantity           & Before update & After update         \\
    \midrule
    Prediction $\yhat$ & $0.637$       & \sol[1.6cm]{$0.643$} \\
    Loss $J$           & $0.132$       & \sol[1.6cm]{$0.127$} \\
    \bottomrule
  \end{tabular}
\end{center}

After one step the prediction has moved closer to the target and the loss
has dropped. The change is small, but repeating it many times is what
training \emph{is}.

% ══════════════════════════════════════════════════════════════════════════════
\section{The Training Loop}

\begin{center}
  \begin{tikzpicture}[
      box/.style={draw=mainblue, fill=boxbg, rounded corners=4pt,
          minimum width=3.8cm, minimum height=1.4cm,
          font=\normalsize, align=center, thick},
      flow/.style={->, thick, mainblue},
      >=Stealth
    ]
    \node[box] (fwd)  at (0,    0)   {Forward pass};
    \node[box] (loss) at (5.2,  0)   {Compute loss\\$J = (y-\yhat)^2$};
    \node[box] (bwd)  at (10.4, 0)   {Backpropagation};
    \node[box] (upd)  at (5.2, -3.0) {Gradient descent\\update};

    \draw[flow] (fwd)  -- node[above, font=\small, black] {$\yhat$} (loss);
    \draw[flow] (loss) -- node[above, font=\small, black] {$J$} (bwd);
    \draw[flow] (bwd.south) |- node[pos=0.25, right, font=\small, black]
      {$\pd{J}{w},\ \pd{J}{b}$} (upd.east);
    \draw[flow] (upd.west) -| node[pos=0.75, left, font=\small, black]
      {new $w,\, b$} (fwd.south);
  \end{tikzpicture}
\end{center}


Check out the interactive demo.
\rimg[0.8]{img/2026-11-10-11-12-30.png}{Interactive demo of one forward--backward pass.}


% ══════════════════════════════════════════════════════════════════════════════
\section{XOR Problem}

In practice we cycle through all four XOR inputs,
average the gradients, and take one update step per batch. After hundreds
or thousands of such cycles the network correctly classifies all four
XOR inputs.

Check out the interactive demo.

\rimg[0.8]{img/2026-11-10-11-14-40.png}{Interactive demo of training on the XOR problem. The network learns to classify the XOR inputs correctly.}
\end{document}
